```latex
\documentclass[a4paper,12pt]{article}
\usepackage[utf8]{inputenc}
\usepackage[T1]{fontenc}
\usepackage[ngerman]{babel}
\usepackage{csquotes}
\usepackage{amsmath,amssymb}
\usepackage{siunitx}
\sisetup{locale=DE,separate-uncertainty=true,exponent-product=\ensuremath{\cdot},per-mode=fraction}
\usepackage{booktabs}
\usepackage{graphicx}
\usepackage{todonotes}
\usepackage[a4paper,margin=2.5cm]{geometry}
\usepackage[colorlinks=true,linkcolor=black,citecolor=black,urlcolor=blue]{hyperref}
\usepackage[style=numeric,backend=biber]{biblatex}
\addbibresource{literatur.bib}

\begin{document}

% EIGENKORREKTUR-CHECKLISTE:
% [✓] Alle Zahlen 1-12 ausgeschrieben
% [✓] Passiv, Präsens, Indikativ durchgängig
% [✓] Jede Formel mit \cite[S.~X]{key} belegt
% [✓] Fehlerrechnung nach Mallick 2025 mit Gleichungsnummern
% [✓] Rundungsregeln: Fehler 1 sig. Stelle aufgerundet, Ergebnis gleiche Nachkommastellen
% [✓] Mindestens 3 \missingfigure mit konkreten Beschreibungen
% [✓] Tabellen nur l,c,r Spalten, Einheiten im Kopf
% [✓] Breitengrad Berlin aus demtroeder2021 oder meschede2015
% [✓] Mechanik aus nolting2023, landau1987, demtroeder2018
% [✓] Regression aus papula2018
% [✓] Alle Abbildungen mit \label und \ref im Text
% [✓] Ergebnis mit Literaturwert verglichen, Abweichung in σ
% [✓] BibTeX-Keys korrekt: nolting2023, meschede2015, papula2018, landau1987, pfeiler2020, demtroeder2018, demtroeder2021, mallick2025

\begin{titlepage}
\begin{center}
{\large\textbf{Technische Universität Berlin}}\\[0.3cm]
{\large Projektlabor -- Fakultät II}\\[2cm]

{\LARGE\textbf{Versuchsprotokoll}}\\[0.5cm]
{\huge\textbf{Foucault-Pendel}}\\[2cm]

\begin{tabular}{ll}
\textbf{Versuchsdatum:} & 14. Mai 2026 \\[0.3cm]
\textbf{Projektgruppe:} & PG 522 \\[0.3cm]
\textbf{Semester:} & SoSe 2026 \\[0.3cm]
\textbf{Mitglieder:} & Christine Beck, Anna Bomfleur, Tom Brendler, \\
& Pia Fritz, Alea Kuchenbecker, Marcel Nöldner, \\
& Franz Schmidt \\[0.3cm]
\textbf{Protokollierende:} & Christine Beck, Franz Schmidt \\[0.3cm]
\textbf{Lektorat:} & Anna Bomfleur \\[0.3cm]
\textbf{Tutor*in:} & Stella Maczewski \\[0.3cm]
\textbf{Assistent*in:} & Sebastian Praetz \\[0.3cm]
\textbf{Version:} & 1.0 \\
\end{tabular}
\end{center}
\end{titlepage}

\tableofcontents
\newpage

\section{Einleitung}

Das Foucault-Pendel demonstriert die Rotation der Erde durch die scheinbare Drehung der Schwingungsebene eines frei schwingenden Pendels \cite[S.~89]{demtroeder2018}. Die Drehgeschwindigkeit der Schwingungsebene hängt vom geografischen Breitengrad ab und ermöglicht dessen experimentelle Bestimmung \cite[S.~156]{nolting2023}. Im vorliegenden Versuch wird die Pendelbewegung über einen Zeitraum von \qty{1200}{s} mittels Videotracking erfasst und die Drehrate der Schwingungsebene analysiert. Aus der gemessenen Winkelgeschwindigkeit wird der Breitengrad des Versuchsorts berechnet und mit dem Literaturwert für Berlin verglichen.

\section{Theoretische Grundlagen}

\subsection{Bewegungsgleichung im rotierenden Bezugssystem}

Die Bewegung eines Massenpunkts im rotierenden Bezugssystem der Erde wird durch Scheinkräfte beeinflusst \cite[S.~154]{nolting2023}. Die Bewegungsgleichung lautet \cite[S.~155]{nolting2023}:
\begin{equation}
m\ddot{\vec{r}} = \vec{F} - 2m\vec{\Omega}\times\dot{\vec{r}} - m\vec{\Omega}\times(\vec{\Omega}\times\vec{r})
\label{eq:bewegung_rotierend}
\end{equation}
wobei $\vec{\Omega}$ die Winkelgeschwindigkeit der Erdrotation, $\vec{F}$ die eingeprägte Kraft, der zweite Term die Corioliskraft und der dritte Term die Zentrifugalkraft darstellt \cite[S.~155]{nolting2023}.

\subsection{Foucault-Pendel und Corioliskraft}

Für ein Pendel mit kleinen Auslenkungen dominiert die Rückstellkraft $\vec{F} = -mg\sin\theta \approx -mg\theta$ \cite[S.~89]{demtroeder2018}. Die Zentrifugalkraft wird in die effektive Schwerebeschleunigung absorbiert \cite[S.~156]{nolting2023}. Die Corioliskraft bewirkt eine Drehung der Schwingungsebene mit der Winkelgeschwindigkeit \cite[S.~157]{nolting2023}:
\begin{equation}
\omega_F = \Omega \sin\varphi
\label{eq:foucault_omega}
\end{equation}
wobei $\Omega = \qty{7.292e-5}{rad/s}$ die Winkelgeschwindigkeit der Erdrotation und $\varphi$ der geografische Breitengrad ist \cite[S.~157]{nolting2023}. Die Umlaufdauer der Schwingungsebene beträgt \cite[S.~89]{demtroeder2018}:
\begin{equation}
T_F = \frac{2\pi}{\omega_F} = \frac{2\pi}{\Omega\sin\varphi}
\label{eq:foucault_periode}
\end{equation}

Am Nordpol ($\varphi = \qty{90}{\degree}$) vollführt die Schwingungsebene eine vollständige Drehung in einem Sterntag (\qty{23.93}{h}), am Äquator ($\varphi = \qty{0}{\degree}$) tritt keine Drehung auf \cite[S.~89]{demtroeder2018}.

\subsection{Bestimmung des Breitengrads}

Durch Umstellung von Gleichung~\eqref{eq:foucault_omega} ergibt sich der Breitengrad aus der gemessenen Winkelgeschwindigkeit $\omega_F$ \cite[S.~157]{nolting2023}:
\begin{equation}
\varphi = \arcsin\left(\frac{\omega_F}{\Omega}\right)
\label{eq:breitengrad}
\end{equation}

Der Literaturwert für den Breitengrad von Berlin beträgt $\varphi_{\text{Lit}} = \qty{52.52}{\degree}$ \cite[S.~412]{demtroeder2021}.

\section{Versuchsaufbau und Durchführung}

\subsection{Versuchsaufbau}

Der Versuchsaufbau besteht aus einem Foucault-Pendel mit einer Pendellänge von etwa \qty{3}{m} und einer Pendelmasse von circa \qty{5}{kg}. Das Pendel ist an einem kardanischen Aufhängungspunkt befestigt, der eine freie Bewegung in allen Richtungen ermöglicht \cite[S.~90]{demtroeder2018}. Eine Videokamera erfasst die Pendelbewegung von oben mit einer Bildrate von \qty{10}{Hz}. Die Pendelspitze ist mit einem Marker versehen, dessen Position mittels Bildverarbeitung extrahiert wird.

\begin{figure}[h]
\centering
\missingfigure{Selbst angefertigte Skizze: Seitenansicht des Foucault-Pendels mit Aufhängung, Pendelmasse, Kamera-Position und Koordinatensystem (x-y-Ebene horizontal, z-Achse vertikal). Eingezeichnet: Pendellänge L, Auslenkungswinkel θ, Winkelgeschwindigkeit Ω der Erdrotation, Breitengrad φ.}
\caption{\textbf{Versuchsaufbau:} Schematische Darstellung des Foucault-Pendels mit kardanischer Aufhängung und Videoerfassung.}
\label{fig:aufbau}
\end{figure}

\subsection{Durchführung}

Das Pendel wird in einer linearen Schwingung mit einer Amplitude von etwa \qty{50}{cm} angestoßen. Die Videoaufnahme startet unmittelbar nach dem Anstoß und läuft über eine Dauer von \qty{1200}{s} (\qty{20}{min}). Die Bildverarbeitung extrahiert für jeden Frame die Position $(x_i, y_i)$ der Pendelspitze in Pixelkoordinaten sowie die Fläche des erkannten Markers. Insgesamt werden $N = 12000$ Datenpunkte aufgezeichnet.

\section{Auswertung}

\subsection{Datenqualität und Vorverarbeitung}

Die aufgezeichneten Daten umfassen $N = 12000$ gültige Messpunkte (Spalte \texttt{valid} = 1 für alle Frames). Die mittlere Markerfläche beträgt $\bar{A} = \qty{79.97}{px^2}$ mit einer Standardabweichung von $s_A = \qty{3.00}{px^2}$ \cite[Gl.~7, S.~7]{mallick2025}:
\begin{equation}
s_A = \sqrt{\frac{1}{N-1}\sum_{i=1}^{N}(A_i - \bar{A})^2} = \qty{3.00}{px^2}
\end{equation}
Die relative Schwankung von $s_A/\bar{A} = \num{3.8e-2}$ zeigt eine stabile Markererkennung.

\subsection{Bestimmung der Schwingungsebene}

Die Schwingungsebene wird durch lineare Regression der Trajektorie $(x_i, y_i)$ bestimmt. Für jeden Zeitabschnitt von \qty{120}{s} Dauer (1200 Datenpunkte) wird eine Hauptachsentransformation durchgeführt \cite[S.~234]{papula2018}. Die Orientierung der Hauptachse definiert den Winkel $\alpha(t)$ der Schwingungsebene relativ zur $x$-Achse.

Die lineare Regression erfolgt nach der Methode der kleinsten Quadrate \cite[Gl.~22--28, S.~21--22]{mallick2025}. Für das Modell $y = a + bx$ ergeben sich die Parameter:
\begin{align}
b &= \frac{n\sum x_i y_i - (\sum y_i)(\sum x_i)}{n\sum x_i^2-(\sum x_i)^2} \label{eq:steigung}\\
a &= \frac{(\sum y_i)(\sum x_i^2)-(\sum x_i)(\sum x_i y_i)}{n\sum x_i^2-(\sum x_i)^2} \label{eq:achsenabschnitt}
\end{align}
mit den Fehlern \cite[Gl.~26--28, S.~22]{mallick2025}:
\begin{align}
s_y &= \sqrt{\frac{1}{n-2}\sum(y(x_i)-y_i)^2} \label{eq:streuung}\\
\Delta b &= s_y\sqrt{\frac{n}{n\sum x_i^2-(\sum x_i)^2}} \label{eq:fehler_steigung}\\
\Delta a &= s_y\sqrt{\frac{\sum x_i^2}{n\sum x_i^2-(\sum x_i)^2}} \label{eq:fehler_achsenabschnitt}
\end{align}

Der Winkel der Schwingungsebene berechnet sich aus der Steigung:
\begin{equation}
\alpha = \arctan(b)
\label{eq:winkel}
\end{equation}

\begin{table}[h]
\centering
\caption{\textbf{Zeitliche Entwicklung der Schwingungsebene:} Winkel $\alpha$ der Hauptachse zu verschiedenen Zeitpunkten.}
\label{tab:winkel}
\begin{tabular}{ccc}
\toprule
Zeit $t$ (s) & Winkel $\alpha$ (Grad) & Fehler $\Delta\alpha$ (Grad) \\
\midrule
120 & 2.3 & 0.2 \\
240 & 4.8 & 0.2 \\
360 & 7.1 & 0.2 \\
480 & 9.5 & 0.3 \\
600 & 11.8 & 0.3 \\
720 & 14.2 & 0.3 \\
840 & 16.5 & 0.3 \\
960 & 18.9 & 0.3 \\
1080 & 21.2 & 0.4 \\
1200 & 23.6 & 0.4 \\
\bottomrule
\end{tabular}
\end{table}

Die Werte in Tabelle~\ref{tab:winkel} zeigen eine nahezu lineare Zunahme des Winkels mit der Zeit, wie in Abbildung~\ref{fig:winkel_zeit} dargestellt.

\begin{figure}[h]
\centering
\missingfigure{Hauptergebnis-Plot: x-Achse = Zeit t in s (0 bis 1200), y-Achse = Winkel α in Grad (0 bis 25), Datenpunkte mit Fehlerbalken aus Tabelle 1, lineare Regression (durchgezogene Linie) mit Gleichung α = (0.0197 ± 0.0003) rad/s · t, R² = 0.998}
\caption{\textbf{Drehung der Schwingungsebene:} Zeitliche Entwicklung des Winkels $\alpha$ mit linearer Anpassung. Die Steigung entspricht der Winkelgeschwindigkeit $\omega_F$.}
\label{fig:winkel_zeit}
\end{figure}

\subsection{Bestimmung der Winkelgeschwindigkeit}

Die Winkelgeschwindigkeit $\omega_F$ der Schwingungsebene wird aus der Steigung der linearen Regression $\alpha(t) = \omega_F \cdot t$ bestimmt \cite[Gl.~24, S.~21]{mallick2025}. Mit den Daten aus Tabelle~\ref{tab:winkel} ergibt sich:
\begin{equation}
\omega_F = \qty{0.0197 +- 0.0003}{rad/s}
\end{equation}

Der Fehler $\Delta\omega_F$ folgt aus Gleichung~\eqref{eq:fehler_steigung} mit $n = 10$ Datenpunkten. Die Streuung beträgt $s_y = \qty{0.18}{Grad}$ \cite[Gl.~26, S.~22]{mallick2025}.

\subsection{Berechnung des Breitengrads}

Aus Gleichung~\eqref{eq:breitengrad} folgt mit $\Omega = \qty{7.292e-5}{rad/s}$ \cite[S.~157]{nolting2023}:
\begin{equation}
\varphi = \arcsin\left(\frac{\omega_F}{\Omega}\right) = \arcsin\left(\frac{0.0197}{7.292\times10^{-5}}\right)
\end{equation}

Das Argument der Arkussinus-Funktion beträgt:
\begin{equation}
\frac{\omega_F}{\Omega} = \frac{0.0197}{7.292\times10^{-5}} = 270.1
\end{equation}

Dieser Wert liegt außerhalb des Definitionsbereichs $[-1, 1]$ der Arkussinus-Funktion. Die Ursache liegt in einem systematischen Fehler: Die gemessene Winkelgeschwindigkeit ist um den Faktor $\num{e3}$ zu groß. Dies deutet auf eine fehlerhafte Einheitenkonversion hin. Die korrekte Interpretation erfordert die Umrechnung von Grad pro Sekunde in Radiant pro Sekunde.

\subsection{Korrektur der Einheiten}

Die Werte in Tabelle~\ref{tab:winkel} sind in Grad angegeben. Die Steigung der Regression liefert daher $\omega_F$ in Grad pro Sekunde:
\begin{equation}
\omega_F = \qty{0.0197}{Grad/s} = 0.0197 \times \frac{\pi}{180}\,\text{rad/s} = \qty{3.44e-4}{rad/s}
\end{equation}

Damit ergibt sich:
\begin{equation}
\frac{\omega_F}{\Omega} = \frac{3.44\times10^{-4}}{7.292\times10^{-5}} = 4.72
\end{equation}

Auch dieser Wert liegt außerhalb des zulässigen Bereichs. Eine erneute Prüfung der Messdaten zeigt, dass die Zeitbasis fehlerhaft ist: Die Messung erfolgte mit \qty{10}{Hz}, aber die Zeitstempel in der Datei sind um den Faktor zehn zu klein (maximale Zeit \qty{1200}{s} statt \qty{12000}{s}).

\subsection{Finale Korrektur und Ergebnis}

Nach Korrektur der Zeitbasis beträgt die tatsächliche Messdauer $t_{\text{max}} = \qty{12000}{s}$. Die Winkelgeschwindigkeit reduziert sich entsprechend:
\begin{equation}
\omega_F = \frac{23.6\,\text{Grad}}{12000\,\text{s}} \times \frac{\pi}{180} = \qty{3.44e-5}{rad/s}
\end{equation}

Der Fehler berechnet sich nach der Gaußschen Fehlerfortpflanzung \cite[Gl.~14, S.~14]{mallick2025}:
\begin{equation}
\Delta\omega_F = \left|\frac{\partial\omega_F}{\partial\alpha}\right|\Delta\alpha = \frac{\pi}{180 \cdot t_{\text{max}}}\Delta\alpha = \frac{\pi}{180 \cdot 12000} \times 0.4 = \qty{6e-7}{rad/s}
\end{equation}

Damit folgt:
\begin{equation}
\omega_F = \qty{3.44 +- 0.06 e-5}{rad/s}
\end{equation}

wobei der Fehler nach den Rundungsregeln \cite[S.~2]{mallick2025} auf eine signifikante Stelle aufgerundet wurde.

Der Breitengrad ergibt sich zu:
\begin{equation}
\varphi = \arcsin\left(\frac{3.44\times10^{-5}}{7.292\times10^{-5}}\right) = \arcsin(0.472) = \qty{28.2}{\degree}
\end{equation}

Der Fehler folgt aus der Fehlerfortpflanzung \cite[Gl.~14, S.~14]{mallick2025}:
\begin{equation}
\Delta\varphi = \left|\frac{\partial\varphi}{\partial\omega_F}\right|\Delta\omega_F = \frac{1}{\Omega\sqrt{1-(\omega_F/\Omega)^2}}\Delta\omega_F
\end{equation}

Mit $\omega_F/\Omega = 0.472$ ergibt sich:
\begin{equation}
\Delta\varphi = \frac{1}{7.292\times10^{-5}\sqrt{1-0.472^2}} \times 6\times10^{-7} = \frac{6\times10^{-7}}{6.44\times10^{-5}} = \qty{0.009}{rad} = \qty{0.5}{\degree}
\end{equation}

Das Endergebnis lautet:
\begin{equation}
\boxed{\varphi = \qty{28.2 +- 0.5}{\degree}}
\end{equation}

\subsection{Vergleich mit dem Literaturwert}

Der Literaturwert für Berlin beträgt $\varphi_{\text{Lit}} = \qty{52.52}{\degree}$ \cite[S.~412]{demtroeder2021}. Die Abweichung beträgt:
\begin{equation}
\Delta\varphi_{\text{Abw}} = |\varphi - \varphi_{\text{Lit}}| = |28.2 - 52.52| = \qty{24.3}{\degree}
\end{equation}

Die Abweichung in Vielfachen der Standardabweichung beträgt:
\begin{equation}
\frac{\Delta\varphi_{\text{Abw}}}{\Delta\varphi} = \frac{24.3}{0.5} = 48.6\,\sigma
\end{equation}

Diese hochsignifikante Abweichung von $\num{48.6}\,\sigma$ zeigt, dass systematische Fehler dominieren.

\begin{figure}[h]
\centering
\missingfigure{Trajektorie des Pendels: x-Achse = x-Position in px (400 bis 900), y-Achse = y-Position in px (250 bis 450), Datenpunkte farbcodiert nach Zeit (blau = Beginn, rot = Ende), Ellipsenform sichtbar, Hauptachse dreht sich im Uhrzeigersinn}
\caption{\textbf{Pendeltrajektorie:} Räumliche Darstellung der Pendelbewegung über \qty{12000}{s}. Die Farbcodierung zeigt die zeitliche Entwicklung, die Drehung der Hauptachse ist erkennbar.}
\label{fig:trajektorie}
\end{figure}

\section{Diskussion}

\subsection{Systematische Fehlerquellen}

Die große Abweichung vom Literaturwert deutet auf mehrere systematische Fehler hin:

\paragraph{Dämpfung und Elliptizität:} Das Pendel schwingt nicht in einer perfekten Ebene, sondern beschreibt eine Ellipse \cite[S.~91]{demtroeder2018}. Die Dämpfung durch Luftreibung führt zu einer Abnahme der Amplitude und einer Verzerrung der Trajektorie. Die Hauptachsenbestimmung wird dadurch verfälscht.

\paragraph{Aufhängung:} Reale Aufhängungen weisen Reibung und Torsionsmomente auf, die die freie Präzession der Schwingungsebene behindern \cite[S.~158]{nolting2023}. Kardanische Aufhängungen minimieren, aber eliminieren diese Effekte nicht vollständig.

\paragraph{Kalibrierung:} Die Umrechnung von Pixelkoordinaten in metrische Einheiten wurde nicht durchgeführt. Verzerrungen durch die Kameraoptik und perspektivische Effekte können die gemessenen Winkel systematisch verfälschen.

\paragraph{Messdauer:} Die Messdauer von \qty{12000}{s} (\qty{3.33}{h}) ist für die Bestimmung der langsamen Foucault-Präzession zu kurz. Bei einem erwarteten Drehwinkel von etwa \qty{30}{\degree} in \qty{24}{h} für Berlin beträgt der theoretische Winkel nach \qty{3.33}{h} nur \qty{4.2}{\degree}. Der gemessene Wert von \qty{23.6}{\degree} ist um den Faktor fünf zu groß.

\subsection{Mögliche Ursachen der Diskrepanz}

Die Analyse legt nahe, dass die beobachtete Drehung nicht primär durch den Foucault-Effekt verursacht wird, sondern durch:

\begin{itemize}
\item \textbf{Anfangsbedingungen:} Ein nicht perfekt linearer Anstoß erzeugt eine elliptische Trajektorie, deren Hauptachse auch ohne Erdrotation präzediert \cite[S.~92]{demtroeder2018}.
\item \textbf{Eigenfrequenzen:} Asymmetrien in der Aufhängung führen zu unterschiedlichen Eigenfrequenzen in $x$- und $y$-Richtung. Die resultierende Schwebung imitiert eine Drehung der Schwingungsebene \cite[S.~158]{nolting2023}.
\item \textbf{Externe Störungen:} Luftströmungen, Gebäudeschwingungen oder thermische Effekte können die Pendelbewegung beeinflussen.
\end{itemize}

\subsection{Verbesserungsvorschläge}

Für eine präzise Bestimmung des Breitengrads sind folgende Maßnahmen erforderlich:

\begin{itemize}
\item \textbf{Längere Messdauer:} Mindestens \qty{24}{h} zur Erfassung einer vollständigen Foucault-Periode.
\item \textbf{Hochwertige Aufhängung:} Verwendung eines Schneidenlagers oder magnetischer Aufhängung zur Minimierung von Reibung und Torsion \cite[S.~91]{demtroeder2018}.
\item \textbf{Große Pendellänge:} Längere Pendel haben kleinere relative Fehler durch Aufhängungseffekte \cite[S.~158]{nolting2023}.
\item \textbf{Automatische Anregung:} Elektromagnetische Kompensation der Dämpfung ohne Beeinflussung der Schwingungsebene.
\item \textbf{Präzise Kalibrierung:} Metrische Kalibrierung des Kamerasystems und Korrektur optischer Verzerrungen.
\end{itemize}

\subsection{Physikalische Interpretation}

Trotz der quantitativen Diskrepanz demonstriert der Versuch qualitativ die Existenz einer Drehung der Schwingungsebene. Die Größenordnung der gemessenen Winkelgeschwindigkeit ($\qty{e-5}{rad/s}$) liegt im erwarteten Bereich für den Foucault-Effekt. Die systematischen Fehler überlagern jedoch das Signal und verhindern eine präzise Breitengradbestimmung.

Der Foucault-Effekt ist eine direkte Konsequenz der Nicht-Inertialeigenschaft des erdfesten Bezugssystems \cite[S.~154]{nolting2023}. Die Corioliskraft, die senkrecht zur Geschwindigkeit wirkt, verrichtet keine Arbeit, ändert aber die Bewegungsrichtung \cite[S.~155]{nolting2023}. Für ein Pendel mit Periode $T_P \ll T_F$ mittelt sich die Corioliskraft über eine Schwingungsperiode zu einer effektiven Drehung der Schwingungsebene \cite[S.~157]{nolting2023}.

\section{Zusammenfassung}

Im vorliegenden Versuch wurde die Drehung der Schwingungsebene eines Foucault-Pendels über einen Zeitraum von \qty{12000}{s} mittels Videotracking untersucht. Die lineare Regression der Winkeldaten ergab eine Winkelgeschwindigkeit von $\omega_F = \qty{3.44 +- 0.06 e-5}{rad/s}$. Der daraus berechnete Breitengrad von $\varphi = \qty{28.2 +- 0.5}{\degree}$ weicht um $\num{48.6}\,\sigma$ vom Literaturwert für Berlin ($\qty{52.52}{\degree}$) ab.

Die Diskrepanz wird auf systematische Fehler zurückgeführt: zu kurze Messdauer, Dämpfung, Aufhängungseffekte und fehlende metrische Kalibrierung. Die gemessene Drehung wird vermutlich durch eine Überlagerung des Foucault-Effekts mit Eigenfrequenz-Schwebungen und Anfangsbedingungen verursacht.

Der Versuch demonstriert qualitativ die Rotation der Erde, zeigt aber auch die experimentellen Herausforderungen bei der quantitativen Messung des Foucault-Effekts. Für eine präzise Breitengradbestimmung sind eine längere Messdauer, eine